% TOP_PROBLEM: {"id": 2851, "title": "Kirby Problem 3.53", "category_id": 7, "category": {"id": 7, "name": "topology", "display_name": "Topology", "description": "Properties preserved under continuous deformations.", "slug": "topology", "order_index": 7, "created_at": "2026-07-31T15:26:25.671Z"}, "status": "open"}
% FIRST_POSTED: null
% ENTRY_KIND: statement_only
% Imported from: batch12.tex; UnsolvedMath ID 2851
\documentclass[11pt]{article}
\usepackage[margin=25mm]{geometry}
\usepackage{fontspec}
\setmainfont{Latin Modern Roman}
\usepackage{amsmath,amssymb,mathtools}
\usepackage{unicode-math}
\setmathfont{Latin Modern Math}
\usepackage{xurl,hyperref}
\hypersetup{colorlinks=true,urlcolor=blue}
\setcounter{secnumdepth}{0}
\setlength{\parindent}{0pt}
\setlength{\parskip}{5pt}
\setlength{\emergencystretch}{3em}
\newcommand{\sref}[2]{\hyperlink{source:#1:#2}{[S#2]}}
\newcommand{\cataloguescope}[1]{\paragraph{Recorded target.}#1}
\begin{document}
% UnsolvedMath ID 2851; source catalogue code KP-3.53
\setcounter{section}{2850}
\section{Kirby Problem 3.53}\label{2851}
{\small\sffamily UnsolvedMath ID 2851 · Topology}
\subsection{Definitions and mathematical statement}

\paragraph{Shared notation.}$\mathbb N=\{1,2,\ldots\}$, and $\mathbb Z,\mathbb Q,\mathbb R,\mathbb C$ denote the integers, rationals, reals and complex numbers. Any other conventions are specified locally.


Write $\mathbb F_2=\mathbb Z/2\mathbb Z$. For a closed oriented $3$-manifold $Y$, $\widehat{\operatorname{HF}}(Y;k)$ denotes hat Heegaard Floer homology with coefficients $k$: it is the homology of $\widehat{\operatorname{CF}}=\operatorname{CF}^-/U\operatorname{CF}^-$, where $\operatorname{CF}^-$ is the Heegaard Floer chain complex over $k[U]$. All spin$^c$ summands are included unless one is specified. A rational homology sphere has $H_*(Y;\mathbb Q)\cong H_*(S^3;\mathbb Q)$. Its first integral homology is finite. An L-space satisfies $\dim_{\mathbb F_2}\widehat{\operatorname{HF}}(Y;\mathbb F_2)=|H_1(Y;\mathbb Z)|$. A pointed genus-$g$ Heegaard diagram is $H=(\Sigma,\boldsymbol\alpha,\boldsymbol\beta,z)$, with two systems of $g$ disjoint attaching circles on $\Sigma$. Its tori $T_\alpha,T_\beta\subset\operatorname{Sym}^g(\Sigma)$ are the images of the products of those circles. The hat chain group is freely generated by $T_\alpha\cap T_\beta$ after making the diagram transverse and admissible. A strong Heegaard diagram for $Y$ has
\[
\#(T_\alpha\cap T_\beta)=|H_1(Y;\mathbb Z)|.
\]
A strong L-space is a rational homology sphere admitting such a diagram. An alternating link has a planar diagram in which each component alternates overcrossings and undercrossings. A link is non-split if no embedded $2$-sphere separates it into two nonempty sublinks. Write $\Sigma_2(L)$ for the double cover of $S^3$ branched along $L$.
\paragraph{Statement recorded in the supplied source.}
For every strong L-space $Y$, does some alternating link $L\subset S^3$ satisfy $Y\cong\Sigma_2(L)$? The link can be taken non-split, as its branched double cover is a rational homology sphere. The assertion ranges over strong diagrams of every genus, not only genus two, and strongness means the generator-count equality above. \sref{2851}{2}
\subsection{Short English statement}
Is every rational homology sphere with a Heegaard diagram having the minimum possible number of Floer generators a double branched cover of an alternating link?
\subsection{Sources}
\begin{description}
\item[\hypertarget{source:2851:1}{S1}] ULAM.AI and the UnsolvedMath Contributors, \emph{UnsolvedMath: A Curated Collection of Open Mathematics Problems}, record 2851 (KP-3.53). CC BY 4.0. \url{https://huggingface.co/datasets/ulamai/UnsolvedMath}.
\item[\hypertarget{source:2851:2}{S2}] R. İnanç Baykur, Robion C. Kirby and Daniel Ruberman (editors), \emph{K3: A New Problem List in Low-Dimensional Topology}, Mathematical Surveys and Monographs 295 (2026), Problem 3.53, p. 169 and following remarks; original author version. \url{https://math.berkeley.edu/sites/default/files/surv-295-ruberman-watermarked-author-pdf.pdf}.
\end{description}
\label{end:2851}
\end{document}
